\documentclass[11pt]{article}

% ====================================================
% PACKAGES
% ====================================================
\usepackage[margin=1in]{geometry}
\usepackage[T1]{fontenc}
\usepackage{lmodern}
\usepackage{microtype}
\usepackage{setspace}
\usepackage{parskip}

\usepackage{amsmath, amssymb, mathtools}
\usepackage{siunitx}

\usepackage{graphicx}
\usepackage{float}
\usepackage{subcaption}

\usepackage{enumitem}
\usepackage{booktabs}

\usepackage{xcolor}
\usepackage{hyperref}

\usepackage{fancyhdr}
\usepackage{lastpage}

\usepackage{listings}
\usepackage{indentfirst}
\usepackage[justification=centering]{caption}

% ====================================================
% PAGE STYLE
% ====================================================
\onehalfspacing%
\setlength{\parindent}{2em}

\pagestyle{fancy}
\fancyhf{}
% Header Design
\lhead{\Assignment\ |\ \course}
\rhead{\Name\ |\ \today}
\cfoot{Page \thepage\ of~\pageref{LastPage}}

\usepackage{listings}
\usepackage{xcolor}

\lstset{
    language=Matlab,
    basicstyle=\ttfamily\small,
    keywordstyle=\color{blue},
    commentstyle=\color{gray},
    stringstyle=\color{orange},
    numbers=left,
    numberstyle=\tiny,
    stepnumber=1,
    numbersep=5pt,
    frame=single,
    breaklines=true,
    showstringspaces=false
}

% ====================================================
% Assignment Info
% ====================================================
\newcommand{\Name}{Arael Anaya}
\newcommand{\Assignment}{Quasi-Static Derivation of Scissor Boom Force}
\newcommand{\course}{Senior Design}


\lstset{
    language=Matlab,
    basicstyle=\ttfamily\small,
    frame=single,
    breaklines=true,
    numbers=left,
    numberstyle=\tiny,
    tabsize=2
}

% makes the hyperlinks look good
\hypersetup{
    colorlinks,
    linkcolor={red!50!black},
    citecolor={blue!50!black},
    urlcolor={blue!80!black}
}


\setlength{\parindent}{1em}
\usepackage{indentfirst}

\begin{document}
\begin{center}
        \includegraphics[width = .8\linewidth]
        {figures/systemDiagram.png}
\end{center}

Assume a \textbf{quasi-static} system, i.e.\ motion is slow and inertia is 
neglected, so the actuator only overcomes friction.

Let $\theta = \text{angle between each scissor link and the horizontal.}$
\[
\begin{aligned}
X(\theta) &= nL\cos\theta, &
\quad & X = \text{horizontal deployed length} \\
Y(\theta) &= nL\sin\theta, &
\quad & Y = \text{vertical actuator position} \\
\frac{dY}{d\theta} &= nL\cos\theta, &
\quad & \text{rate of vertical motion w.r.t. }\theta
\end{aligned}
\]

Resistance consists of hinge and rail friction. With hinge torque $\tau_h$ and $n_h$ hinges per bay,
\[
\tau_{\text{hinges}} = n n_h \tau_h, 
\qquad
\tau_{\text{rail}} = F_{\text{rail}} \frac{dY}{d\theta}.
\]
Thus,
\[
\tau_\theta = \tau_{\text{hinges}} + \tau_{\text{rail}}.
\]

Using virtual work\footnote{Since $Y$ and $\theta$ are in different
coordinate systems, we must use virtual work and more precisely the principle
of least action applied to energy in a static system.},
\[
F_{\text{act}}\,dY = \tau_\theta\,d\theta
\;\;\Rightarrow\;\;
F_{\text{act}} = \frac{\tau_\theta}{dY/d\theta}.
\]

\[
\boxed{
F_{\text{act}}(\theta)
=
\frac{
n n_h \tau_h + F_{\text{rail}} (nL\cos\theta)
}{
nL\cos\theta
}
}
\]
This is the quasi-static actuator force required to move the scissor boom 
at angle \(\theta\).

\newpage

%%%%%%%%%%%%%%%%%%%%%%%%%%%%%%%%%%%%%%%%%%%%%%%%%%%%

% \begin{center}
%         \includegraphics[width = .3\linewidth]
%         {figures/jointDiagram.png}
% \end{center}

Let $\alpha_k = \text{absolute angle of link } k,$ and
$\theta_i = \alpha_{i+1} - \alpha_i$, where $\theta = 0$ is a fully straight joint.

\[
\alpha_k = \alpha_k(t) + \sum_{j=1}^{k-1} \theta_j,
\quad k = 2,\dots,n
\]

\[
\begin{aligned}
X &= \sum_{k=1}^{n} L \cos(\alpha_k),
&\quad & X = \text{horizontal tip position} \\[6pt]
Y &= \sum_{k=1}^{n} L \sin(\alpha_k),
&\quad & Y = \text{vertical tip position}
\end{aligned}
\]

Initial folded configuration:
\[
\begin{aligned}
\alpha_k(0) &= \frac{\pi}{2} + (-1)^{k-1}\phi_0,
&\quad & \text{alternating link angles} \\[6pt]
\theta_i(0) &= \pm 2\phi_0,
&\quad & \text{alternating joint angles}
\end{aligned}
\]

Model each compliant joint as a circular arc:
\[
R = \frac{s}{\theta},
\qquad
\mathbf{r} =
\begin{bmatrix}
R \sin\theta \\
R(1-\cos\theta)
\end{bmatrix}
\qquad
\mathbf{F}_p = pA
\begin{bmatrix}
\cos\theta \\
\sin\theta
\end{bmatrix}
\]


Moment about the joint base:
\[
M_p = \mathbf{r} \times \mathbf{F}_p \to \boxed{M_p = p A \frac{s}{\theta} (1 - \cos\theta)
}
\]

Joint resistance is modeled as a torsional spring:
\[
\boxed{
M_r = k \theta
}
\]

Quasi-static equilibrium:
\[
M_p = M_r \to p A \frac{s}{\theta}(1 - \cos\theta) = k \theta
\]

\[
\boxed{
p(\theta) =
\frac{k \theta^2}{A s (1 - \cos\theta)}
}
\]

This is the quasi-static pressure required to actuate each compliant joint 
as a function of $\theta$.


\end{document}